\documentclass[a4paper,UKenglish,cleveref,thm-restate]{lipics-v2021}

\bibliographystyle{plainurl}% the mandatory bibstyle

\title{Inception algebras / jumps with data-passing} %TODO

\author%
{Vikraman Choudhury}%
{University of Strathclyde, Glasgow, United Kingdom}%
{vikraman.choudhury@strath.ac.uk}%
{https://orcid.org/0000-0003-2030-8056}%
{Supported by ...}%TODO

\author%
{Marcelo Fiore}%
{University of Cambridge, Cambridge, United Kingdom}%
{marcelo.fiore@cst.cam.ac.uk}%
{...}%TODO
{Supported by ...}%TODO

\author%
{Gregor Feierabend}%
{University of Cambridge, Cambridge, United Kingdom}%
{gregor.feierabend@cst.cam.ac.uk}%
{...}%TODO
{Supported by ...}%TODO

\authorrunning{V. Choudhury, M. Fiore and G. Feierabend} %TODO mandatory. First: Use abbreviated first/middle names. Second (only in severe cases): Use first author plus 'et al.'

\Copyright{Vikraman Choudhury, Marcelo Fiore and Gregor Feierabend} %TODO mandatory, please use full first names. LIPIcs license is "CC-BY";  http://creativecommons.org/licenses/by/3.0/

\ccsdesc[100]{\textcolor{red}{Replace ccsdesc macro with valid one}} %TODO mandatory: Please choose ACM 2012 classifications from https://dl.acm.org/ccs/ccs_flat.cfm 

\keywords{computational effects, continuations} %TODO mandatory; please add comma-separated list of keywords

\acknowledgements{We want to thank \dots}%optional

\RequirePackage{currfile}
\providecommand{\ifarxiv}[2]{#2}
\ifcurrfilename{arxiv-inc.tex}{%
  \renewcommand{\ifarxiv}[2]{#1}%
}{}

\ifarxiv{%
\pdfoutput=1
\hideLIPIcs
\nolinenumbers
\usepackage[appendix=inline,bibliography=common]{apxproof}
}{%
\usepackage[appendix=strip,bibliography=common]{apxproof}
}

\usepackage{xspace}
\usepackage{longtable}
\usepackage[automake,xindy={language=english},symbols,abbreviations,sort=def]{glossaries-extra}
%%%%%%%%%%%%%%%%%%%%%%%%%%%%%%%%%%%%%%%%%%%%%%%%%%%%%%%%%%%%%%%%%%%%%%%%%%
% Symbol definitions
%%%%%%%%%%%%%%%%%%%%%%%%%%%%%%%%%%%%%%%%%%%%%%%%%%%%%%%%%%%%%%%%%%%%%%%%%%

\glsxtrnewsymbol[description={Category of contexts (\textsf{FinSet})}]{CatOfCtx}{\ensuremath{\mathbb{F}}}

\glsxtrnewsymbol[description={Continuation Monad $R^{R^{-}}$ on ${\textsf{Set}}^{\glspt{CatOfCtx}}$ with $R \in {\textsf{Set}}^{\glspt{CatOfCtx}}$}]{ContMonad}{\ensuremath{K_R}}

\glsxtrnewsymbol[description={Variable presheaf $Y(1) \in {\textsf{Set}}^{\glspt{CatOfCtx}}$}]{VarPresheaf}{\ensuremath{L}}

\glsxtrnewsymbol[description={Label type}]{LabelType}{\ensuremath{\ell}}

\glsxtrnewsymbol[description={Context},user1={\ensuremath{\Gamma}},user2={\ensuremath{\Delta}},user3={\ensuremath{\Psi}}]{Ctx}{\ensuremath{\Gamma, \Delta, \Psi}}

\glsxtrnewsymbol[description={Pure term},user1={\ensuremath{V}},user2={\ensuremath{W}}]{Pure}{\ensuremath{V, W}}

\glsxtrnewsymbol[description={Effectful term},user1={\ensuremath{M}},user2={\ensuremath{N}}]{Comp}{\ensuremath{M, N}}

\glsxtrnewsymbol[description={Variable},user1={\ensuremath{x}},user2={\ensuremath{y}},user3={\ensuremath{z}}]{Var}{\ensuremath{x, y, z}}

\glsxtrnewsymbol[description={Unit term}]{UnitTerm}{\ensuremath{\ast}}

\glsxtrnewsymbol[description={Lambda abstraction operator}]{LamOp}{\textsf{lam}}

\glsxtrnewsymbol[description={Pair operator}]{PairOp}{\textsf{pair}}

\glsxtrnewsymbol[description={Push operator}]{PushOp}{\textsf{push}}

\glsxtrnewsymbol[description={Return operator}]{ReturnOp}{\textsf{return}}

\glsxtrnewsymbol[description={Application operator}]{AppOp}{\textsf{app}}

\glsxtrnewsymbol[description={Jump operator}]{VarOp}{\textsf{var}}

\glsxtrnewsymbol[description={Branch operator}]{SubOp}{\textsf{sub}}

\glsxtrnewsymbol[description={Pattern match operator}]{PmOp}{\textsf{pm}}

\glsxtrnewsymbol[description={Type},user1={\ensuremath{X}},user2={\ensuremath{Y}},user3={\ensuremath{Z}}]{Type}{\ensuremath{X, Y, Z}}

\glsxtrnewsymbol[description={Membership judgement},user1={\ensuremath{\in}}]{Member}{\ensuremath{x : X \in \Gamma}}

\glsxtrnewsymbol[description={Typing judgement for pure terms},user1={\ensuremath{\vdash^v}}]{PureJudge}{\ensuremath{\Gamma \vdash^v V : X}}

\glsxtrnewsymbol[description={Typing judgement for computations},user1={\ensuremath{\vdash^c}}]{CompJudge}{\ensuremath{\Gamma \vdash^c M : X}}

\glsxtrnewsymbol[description={Appending a type to a context},user1={\ensuremath{\cdot}}]{CtxCons}{\ensuremath{\Gamma \cdot X}}

\glsxtrnewsymbol[description={Function type},user1={\ensuremath{\Rightarrow}}]{ArrowType}{\ensuremath{X \Rightarrow Y}}

\glsxtrnewsymbol[description={Unit type}]{UnitType}{\ensuremath{\mathbf{1}}}

\glsxtrnewsymbol[description={Product type},user1={\ensuremath{\times}}]{ProdType}{\ensuremath{X \times Y}}

%%%%%%%%%%%%%%%%%%%%%%%%%%%%%%%%%%%%%%%%%%%%%%%%%%%%%%%%%%%%%%%%%%%%%%%%%%
% Macros
%%%%%%%%%%%%%%%%%%%%%%%%%%%%%%%%%%%%%%%%%%%%%%%%%%%%%%%%%%%%%%%%%%%%%%%%%%

\newcommand\glsi[1]{\glshyperlink[\glsentryuseri{#1}]{#1}\glsadd{#1}\xspace}
\newcommand\glsii[1]{\glshyperlink[\glsentryuserii{#1}]{#1}\glsadd{#1}\xspace}
\newcommand\glsiii[1]{\glshyperlink[\glsentryuseriii{#1}]{#1}\glsadd{#1}\xspace}

%%%%%%%%%%%%%%%%%%%%%%%%%%%%%%%%%%%%%%%%%%%%%%%%%%%%%%%%%%%%%%%%%%%%%%%%%%
% Custom glossary style
%%%%%%%%%%%%%%%%%%%%%%%%%%%%%%%%%%%%%%%%%%%%%%%%%%%%%%%%%%%%%%%%%%%%%%%%%%

\newlength{\myglosssymbolwidth}\setlength{\myglosssymbolwidth}{0.15\textwidth}
\newlength{\myglossdescwidth}\setlength{\myglossdescwidth}{0.55\textwidth}
\newlength{\myglosskeywidth}\setlength{\myglosskeywidth}{0.2\textwidth}

\newglossarystyle{long-symbols}{%
% put the glossary in a longtable environment:
\renewenvironment{theglossary}%
%% To determine width of colums automatically: {\begin{longtable}{lp{\glsdescwidth}p{\glspagelistwidth}}}%
{\begin{longtable}{p{\myglosssymbolwidth}p{\myglossdescwidth}p{\myglosskeywidth}}}%
{\end{longtable}}%
% Set the table's header:
\renewcommand*{\glossaryheader}{%
\bfseries Symbol & \bfseries Description & \bfseries Key
\\\endhead}%
% No heading between groups:
\renewcommand*{\glsgroupheading}[1]{}%
\renewcommand*{\glossentry}[2]{%
\glstarget{##1}{\glossentryname{##1}}
& \glossentrydesc{##1}% Description
& \texttt{##1}%
\tabularnewline % end of row
}%
% Similarly for sub-entries (no sub-entry numbers)
\renewcommand*{\subglossentry}[3]{%
% ignoring first argument (sub-level)
\glstarget{##2}{\glossentryname{##2}}
& \glossentrydesc{##2}% Description
& \texttt{##2}%
\tabularnewline % end of row
}%
% Nothing between groups:
\renewcommand*{\glsgroupskip}{}%
}

\makeglossaries
\overfullrule=5pt

\usepackage{mathpartir}

%Editor-only macros:: begin (do not touch as author)%%%%%%%%%%%%%%%%%%%%%%%%%%%%%%%%%%
\EventEditors{John Q. Open and Joan R. Access}
\EventNoEds{2}
\EventLongTitle{42nd Conference on Very Important Topics (CVIT 2016)}
\EventShortTitle{CVIT 2016}
\EventAcronym{CVIT}
\EventYear{2016}
\EventDate{December 24--27, 2016}
\EventLocation{Little Whinging, United Kingdom}
\EventLogo{}
\SeriesVolume{42}
\ArticleNo{23}
%%%%%%%%%%%%%%%%%%%%%%%%%%%%%%%%%%%%%%%%%%%%%%%%%%%%%%

\begin{document}

\maketitle

%TODO mandatory: add short abstract of the document
\begin{abstract}
Lorem ipsum dolor sit amet, consectetur adipiscing elit. Praesent convallis orci arcu, eu mollis dolor. Aliquam eleifend suscipit lacinia. Maecenas quam mi, porta ut lacinia sed, convallis ac dui. Lorem ipsum dolor sit amet, consectetur adipiscing elit. Suspendisse potenti. 
\end{abstract}

\section{For the authors: Macros and Notation}
I have added a glossary in \verb+glossary.tex+.
This is not supposed to be included in the LICS version, but we could maybe include it in the Arxiv version.
The provisional glossary below shows all the defined entries and their corresponding keys (just to make things easier for us while writing).
A symbol with a corresponding key \texttt{somekey} can be printed by \verb+\gls{somekey}+ in the document.

\section{definitions}

\begin{figure}[t]

\begin{subfigure}{0.95\linewidth}
  \begin{align*}
    \text{Pure terms:} && \glsi{Pure}, \glsii{Pure} &:=\; \glsi{Var} \mid \gls{UnitTerm} \mid \gls{LamOp} ( \glsi{Comp} ) \mid \gls{PairOp} (\glsi{Pure} , \glsii{Pure})
    \\
    \text{Effectful terms / computations:} && \glsi{Comp}, \glsii{Comp} &:=\; \gls{ReturnOp} (\glsi{Pure}) \mid \gls{PushOp} ({\glsi{Comp}} , {\glsii{Comp}}) \mid \gls{AppOp} (\glsi{Pure} , \glsii{Pure}) \mid \\
         &&& \quad\quad \gls{VarOp} (\glsi{Pure}) \mid \gls{SubOp} (\glsi{Comp}, \glsii{Comp}) \mid \gls{PmOp} ( \glsi{Pure} , \glsi{Comp} )
  \end{align*}
  \caption{Grammar for terms}
  \label{fig:terms-a}
\end{subfigure}

\begin{subfigure}{0.95\linewidth}
  \begin{mathpar}
    \inferrule
      {\glsi{Var} : \glsi{Type} \mathbin{\glsi{Member}} \glsi{Ctx}}
      {\glsi{Ctx} \mathbin{\glsi{PureJudge}} \glsi{Var} : \glsi{Type}}

    \inferrule
      {\glsi{Ctx} \mathbin{\glsi{CtxCons}} \glsi{Type} \mathbin{\glsi{CompJudge}} \glsi{Comp} : \glsii{Type}}
      {\glsi{Ctx} \mathbin{\glsi{PureJudge}} \gls{LamOp} ( \glsi{Comp} ) : \glsi{Type} \mathbin{\glsi{ArrowType}} \glsii{Type}}

    \inferrule
      { }
      {\glsi{Ctx} \mathbin{\glsi{PureJudge}} \gls{UnitTerm} : \gls{UnitType}}

    \inferrule
      {\glsi{Ctx} \mathbin{\glsi{PureJudge}} \glsi{Pure} : \glsi{Type} \\ \glsi{Ctx} \mathbin{\glsi{PureJudge}} \glsii{Pure} : \glsii{Type}}
      {\glsi{Ctx} \mathbin{\glsi{PureJudge}} \gls{PairOp} ( \glsi{Pure} , \glsii{Pure} ) : \glsi{Type} \mathbin{\glsi{ProdType}} \glsii{Type}}

  \end{mathpar}
  \caption{Typing rules for pure terms}
  \label{fig:terms-b}
\end{subfigure}

\begin{subfigure}{0.95\linewidth}
  \begin{mathpar}
    \inferrule
      {\glsi{Ctx} \mathbin{\glsi{PureJudge}} \glsi{Pure} : \glsi{Type}}
      {\glsi{Ctx} \mathbin{\glsi{CompJudge}} \gls{ReturnOp} (\glsi{Pure}) : \glsi{Type}}

    \inferrule
      {\glsi{Ctx} \mathbin{\glsi{CompJudge}} \glsi{Comp} : \glsi{Type} \\ \glsi{Ctx} \mathbin{\glsi{CtxCons}} \glsi{Type} \mathbin{\glsi{CompJudge}} \glsii{Comp} : \glsii{Type}}
      {\glsi{Ctx} \mathbin{\glsi{CompJudge}} \gls{PushOp} (\glsi{Comp} , \glsii{Comp}) : \glsii{Type}}

    \inferrule
      {\glsi{Ctx} \mathbin{\glsi{PureJudge}} \glsi{Pure} : \glsi{Type} \mathbin{\glsi{ArrowType}} \glsii{Type} \\ \glsi{Ctx} \mathbin{\glsi{PureJudge}} \glsii{Pure} : \glsi{Type}}
      {\glsi{Ctx} \mathbin{\glsi{CompJudge}} \gls{AppOp}(\glsi{Pure}, \glsii{Pure}) : \glsii{Type}}

    \inferrule
      {\glsi{Ctx} \mathbin{\glsi{PureJudge}} \glsi{Comp} : \gls{LabelType}}
      {\glsi{Ctx} \mathbin{\glsi{CompJudge}} {\gls{VarOp}}_{\glsi{Type}} (\glsi{Comp}) : \glsi{Type}}

    \inferrule
      {\glsi{Ctx} \mathbin{\glsi{CtxCons}} \gls{LabelType} \mathbin{\glsi{CompJudge}} \glsi{Comp} : \glsi{Type} \\ \glsi{Ctx} \mathbin{\glsi{CompJudge}} \glsii{Comp} : \glsi{Type}}
      {\glsi{Ctx} \mathbin{\glsi{CompJudge}} {\gls{SubOp}}_{\glsi{Type}} (\glsi{Comp} , \glsii{Comp}) : \glsi{Type}}

    \inferrule
      {\glsi{Ctx} \mathbin{\glsi{PureJudge}} \glsi{Pure} : \glsi{Type} \mathop{\glsi{ProdType}} \glsii{Type} \\ \glsi{Ctx} \mathbin{\glsi{CtxCons}} \glsi{Type} \mathbin{\glsi{CtxCons}} \glsii{Type} \mathbin{\glsi{CompJudge}} \glsi{Comp} : \glsiii{Type}}
      {\glsi{Ctx} \mathbin{\glsi{CompJudge}} \gls{PmOp}(\glsi{Pure}, \glsi{Comp}) : \glsiii{Type}}
  \end{mathpar}
  \caption{Typing rules for effectful terms / computations}
  \label{fig:terms-c}
\end{subfigure}

\caption{Terms}
\label{fig:terms}
\end{figure}

\clearpage
\glsaddall
\printglossary[type=symbols, title={Glossary}, style={long-symbols}]

\end{document}
